\newcommand{\sedarchcaption}{
\textbf{Model architecture} \\
\small
(\textbf{a}) Natural, ``real-world'' features $x$ are encoded by neural activity $y$. In this example, three active (of six total) features are simultaneously represented by the joint activity of three neurons.
(\textbf{b}) An SED reconstructs target neural activity $\psi$ as $\hat{\psi}$ from input $y$ through sparsely active latents $z$. When training is successful, the learned latents $z$ correspond to natural features $x$. Reconstructing the input population itself ($\psi = y$) gives an autoencoder; reconstructing a separate but related target population gives a transcoder. By default, one training example is a single timebin.
(\textbf{c}) A Matryoshka SED divides the latents into nested levels $L_1 \subset L_2 \subset L_3$, each of which attempts a full reconstruction of $\psi$. Black dashed boxes mark the latents included at a given level and orange dashed boxes mark latents added at higher levels. Here, batch top-$k$ selects one additional active latent at each level (yellow), allowing lower (higher) levels to capture coarser (finer) features.
(\textbf{d}) A W-SED instead treats a sequence of $S$ timebins of $N$ recorded units, $\mathbf{Y}\in\mathbb{R}^{S\times N}$, as one training example and reconstructs all $S$ bins of a target $\mathbf{\Psi}\in\mathbb{R}^{S\times N}$ as $\hat{\mathbf{\Psi}}$ (the batch dimension is omitted throughout). A FW-SED encoder linearly projects the vectorized window to the $D$ latents $\mathbf{z}$; alternatively, a TW-SED encoder treats each timebin's population state as a token, integrates temporal context via self-attention, and projects the temporally contextualized final token to the latents. The latents pass through ReLU and batch top-$k$ (optionally with Matryoshka levels), and each active latent contributes a timebin-by-unit activity pattern to the full reconstruction $\hat{\mathbf{\Psi}}$. The loss is a weighted average of per-timebin losses, $\sum_s w_s \ell_s / \sum_s w_s$; reconstruction quality is also reported per timebin. An optional shift-equivariant mode (each latent is a motif that may occur at any timebin) is described in \hyperref[subsubsection:model_architecture_training]{\ref*{subsubsection:model_architecture_training}}.
}